\section{Problem Definition}
\label{sec:problem}

\subsection{Routing setting}
\label{sec:routing-setting}

We consider a fixed \textbf{model pool} $\Mpool=\{M_1,\ldots,M_K\}$ whose members have \textbf{distinct, documented capabilities}---they are not interchangeable.
At Stage~3 we designate a \textbf{primary pair} for binary routing: $\Mlo$ (lower capability, lower cost) and $\Mhi$ (higher capability, higher cost), with $\mathrm{cap}(\Mlo)<\mathrm{cap}(\Mhi)$.
The pool may contain additional members for comparative signals; the primary pair defines oracle $r(q)$ and deployable policy $\pi(q)$.
In this paper $|\mathcal{P}|=2$, so \textbf{binary LLM routing} reduces to selecting $\Mlo$ or $\Mhi$ per query.

For each query $q$, pool members run under a frozen prompting and decoding protocol; we write $y(q,M)\in\{0,1\}$ for correctness under an \textbf{objective} task metric (Section~\ref{sec:experimental-setup}).
Let $y_{\mathrm{lo}}=y(q,\Mlo)$ and $y_{\mathrm{hi}}=y(q,\Mhi)$.

Our goal is \textbf{appropriate model selection}: from signals $x(q)$, choose the pool member that balances correctness and cost.
We evaluate \textbf{binary model-selection} policies by task accuracy and average model cost on a held-out \textbf{test} split; an offline \textbf{oracle} upper bound is defined below for evaluation only.
We do \textbf{not} propose a new cascade algorithm---sequential multi-model invoke appears only in related work (Section~\ref{sec:related}).

\subsection{Routing need (oracle label)}
\label{sec:routing-need}

To analyze and fit routers on calib, we need an offline label for whether $\Mhi$ was the \textbf{appropriate} pool member ex post.
We define \textbf{routing need} $r(q)$---also \textbf{escalation utility} in our notation---as an \textbf{oracle label}, not the deployable policy $\pi(q)$:
\begin{equation}
  r(q) = \mathbb{1}\big[\, y(q,\Mlo)=0 \;\wedge\; y(q,\Mhi)=1 \,\big].
  \label{eq:routing-need}
\end{equation}
$r(q)=1$ iff $\Mlo$ failed but $\Mhi$ would succeed---i.e.\ $\Mhi$ was the right choice, not $\Mlo$.
This is \textbf{not} generic task difficulty: \emph{too\_hard} queries ($y_{\mathrm{lo}}=y_{\mathrm{hi}}=0$) have $r(q)=0$ because $\Mhi$ does not improve the outcome.
Queries fall into four buckets according to $(y_{\mathrm{lo}},y_{\mathrm{hi}})$:

\begin{table}[t]
  \centering
  \small
  \begin{tabular}{@{}llcl@{}}
    \toprule
    \textbf{Bucket} & $(y_{\mathrm{lo}}, y_{\mathrm{hi}})$ & $r(q)$ & \textbf{Appropriate model (oracle)} \\
    \midrule
    easy         & $(1,1)$ & 0 & $\Mlo$ \\
    opportunity  & $(0,1)$ & \textbf{1} & $\Mhi$ \\
    lo\_only     & $(1,0)$ & 0 & $\Mlo$ \\
    too\_hard    & $(0,0)$ & 0 & $\Mlo$ \\
    \bottomrule
  \end{tabular}
  \caption{Query buckets and oracle $r(q)$. \emph{Opportunity} queries have $r(q)=1$.}
  \label{tab:buckets}
\end{table}

We also write the correctness gap $\Delta(q)=y_{\mathrm{hi}}-y_{\mathrm{lo}}$; $r(q)=1$ corresponds to $\Delta(q)=+1$.
For \textbf{evaluation only}, the \textbf{oracle} policy $\pi^*(q)$ selects the cheapest correct pool member: $\Mlo$ if $y_{\mathrm{lo}}=1$; else $\Mhi$ if $y_{\mathrm{hi}}=1$; else $\Mlo$.
The oracle uses offline labels and is not a deployable router.

\subsection{Oracle label vs.\ deployable policy}
\label{sec:oracle-vs-policy}

\begin{table}[t]
  \centering
  \small
  \begin{tabular}{@{}lll@{}}
    \toprule
    & \textbf{Oracle label $r(q)$} & \textbf{Deployable policy $\pi(q)$} \\
    \midrule
    When known & Offline (both models on calib/test) & Online at routing time \\
    Inputs & $y_{\mathrm{lo}}, y_{\mathrm{hi}}$ & Signals $x(q)$ only \\
    Question & Was $\Mhi$ the appropriate choice ex post? & Which pool member should answer? \\
    \bottomrule
  \end{tabular}
  \caption{Oracle supervision vs.\ binary model selection.}
  \label{tab:oracle-vs-policy}
\end{table}

Deployable routing: compute $x(q)$, score $s(q)$, and \textbf{select} $\Mlo$ or $\Mhi$ via $\pi(q)$ (Section~\ref{sec:routing-policy}).
$\pi(q)$ never observes $r(q)$, $y_{\mathrm{lo}}$, or $y_{\mathrm{hi}}$ at decision time.

Routing need $r(q)$ is computed offline for analysis and policy fit; it is \textbf{never} an input to signal extraction (Section~\ref{sec:signal-extraction}).

\subsection{What ``unsupervised'' means}
\label{sec:unsupervised-def}

\textbf{Unsupervised} describes \textbf{how signals are obtained} (layer~1), not the absence of all labels in the paper.
We use three \textbf{supervision layers}:

\begin{table}[t]
  \centering
  \small
  \begin{tabular}{@{}clcc@{}}
    \toprule
    \textbf{Layer} & \textbf{Program §} & Uses $r(q)$? & \textbf{Unsupervised?} \\
    \midrule
    1 --- Signal extraction   & Section~\ref{sec:signal-extraction} & No  & Yes \\
    2 --- Signal analysis     & Section~\ref{sec:signal-analysis}   & Yes, calib only & Analysis supervised; signals not \\
    3 --- Policy fit          & Section~\ref{sec:routing-policy}      & Yes, calib only & Combination supervised; definitions not \\
    \bottomrule
  \end{tabular}
  \caption{Three supervision layers. ``Unsupervised signals'' refers to layer~1 only.}
  \label{tab:layers}
\end{table}

At layer~1, each feature $x_i(q)$ is computed from query text and/or model responses only---no routing preferences, no routing classifier training data.
\textbf{Model-response} signals require running pool members on $q$; ``unsupervised'' refers to \textbf{routing supervision}, not compute.
At layers~2--3, we use offline $r(q)$ on \textbf{calib} to score informativeness (H1--H3) and to fit policy parameters $(\lambda,\tau)$ (H4); all reported routing metrics use locked parameters on \textbf{test} only.

Combining signals with calib labels follows the advisor formulation: unsupervised \emph{features}, supervised \emph{weights}.
This differs from supervised routers that learn routing inputs from preference datasets end-to-end, and from claims of ``zero labels'' in the entire system.

We do \textbf{not} mean classical unsupervised routing (e.g., clustering queries with no oracle on calib) as our method, nor every paper that uses the word ``unsupervised'' for synthetic judge scores in retrieval pipelines.

\subsection{Research questions}

\begin{enumerate}
  \item[RQ1.] What \textbf{model-independent} (query-derived), \textbf{model-dependent} (model-response), and \textbf{cross-model} (cross-model comparative) signals can we estimate without routing labels at extraction time?
  \item[RQ2.] Which layers predict oracle $r(q)$ (appropriate model choice) on calib? \textbf{(H1--H3)}
  \item[RQ3.] After freezing $x(q)$, do learned weights outperform hand-designed rules on the same features? \textbf{(H4)}
\end{enumerate}

Signal analysis (RQ2) and routing evaluation (RQ3) are \textbf{separate} claims: informative signals can fail to improve Pareto routing, and good routing performance can arise primarily from policy fit rather than signal structure.
Routing quality is judged by $\mathrm{Acc}(\pi)$ and $\mathrm{Cost}(\pi)$ on test with \textbf{Pareto} dominance as the primary comparison (Section~\ref{sec:experimental-setup}).
The routing score $s(q)$ and policy $\pi(q)$ are defined once in Section~\ref{sec:routing-policy}.

\subsection{Scope and universality}
\label{sec:scope}

Our \textbf{methodology} is setting-agnostic; our \textbf{empirical results} are conditional on one locked Experimental Setting $\mathcal{S}=(\mathcal{D},\mathcal{P},\Pi)$.

\paragraph{Universal (method).}
Oracle $r(q)$ for appropriate-model analysis; deployable \textbf{binary model selection} $\pi(q)$ from pre-specified signals; three information layers ($\phi$, $\psi$, $\chi$); parallel H1--H3; freeze $x(q)$; Pareto routing (H4).
These apply wherever pool outputs are objectively gradable (MCQ match, unit tests, grader scores).

\paragraph{Setting-specific (instance).}
Benchmark, pool pair, prompt protocol, and concrete extractors (e.g.\ choice-letter entropy on ARC; load/ambiguity descriptors for $\phi$).
The same layers apply to open-ended tasks with different $\psi$ extractors---we do not claim ARC features transfer without refit.

\paragraph{Relation to routing surveys.}
Surveys distinguish \textbf{routing} (choose a model) from \textbf{cascading} (sequential invoke).
Our contribution is the former: signals $\to$ $\pi(q)$ $\to$ one pool member.
FrugalGPT-style cascades are related work only (Section~\ref{sec:related}).
Surveys also classify routers by \emph{when} the decision is made, \emph{what information} is available, and \emph{how} the score is computed; our layers map to information availability (Section~\ref{sec:intro}).
